\section{22 September 2026: A common pre-adoption source for historical proposals}

Historical proposals filed in 2019--2022 now use the inactive-inclusive
23Q4 Housing Database. Jacob selected all recorded proposals with status
flags. The archived unit count, filing date, lot, building identifier,
description and ownership category define the historical record. The latest
DOB update in this archive is 12 January 2024. It measures designs recorded
before 485-x adoption, including amendments made since their initial filings.
Post-policy proposals retain the 25Q4 Housing Database count with DOB fallback
when that count is missing. Recorded zeros remain zero.

The source correction gives Livingston one parent with 144+105=249 units.
Both archived filings identify the same lot and refer to the neighboring
application. The later 99-unit design does not enter its historical outcome.
Across the historical panel, exact-99 constituent filings fall from six to
two. Four of the previous six have other unit counts in 23Q4. This comparison
identifies a source-vintage problem; it does not establish when those designs
changed or why.

\paragraph{Inactive proposals and alternatives.}
Archived status and membership in the archive's active/completed file are
preserved separately. Earlier withdrawn applications count as alternatives
when a valid archived BIN, identical lot and exact address identify one later
nonwithdrawn application within the parent window. Multiple withdrawn
versions can share that successor. In the canonical historical cohorts this
rule excludes 34 earlier applications carrying 734 units from additive totals.
It preserves the initial date, successor filing date and evidence basis;
the archive does not supply the date of withdrawal.

Two reviewed roles supplement that rule. The 773 Neptune institutional job
retains its recorded 189 Class~A units as source evidence but contributes no
residential units: its archived description specifies a house of worship and
the later official DOB filing records institutional use with zero dwellings.
The withdrawn 158-unit West 49th Street proposal is an alternative to the
158-unit West 48th Street application. Their archived lot and program agree,
and the City's environmental assessment describes one housing building on
the reviewed western ground. The assessment does not name both application
IDs, so this relationship remains a documented inference. Its earlier date
anchors the parent; the reviewed 21,924-square-foot allocation is unchanged.

The panel contains 2,173 historical parents and 2,348 additive constituents,
with 141,632 units, compared with 1,811 parents, 1,904 constituents and
129,335 units before the source correction. Of the new parents, 206 include
an inactive additive constituent. Another 25 have active additive
constituents but an inactive earlier application. Separate parent flags
record these two definitions of activity. The post-policy panel retains
907 parents, 1,005 constituents and 61,319 units; the comparison also verifies
its dates, memberships, exposure, composition eligibility and land covariates.

The underlying 2019--2022 archive contains 2,432 new-building filings with at
least six units: 2,166 active/completed-file records and 266 other records.
The existing positive-area pre-filing parcel-coverage rule leaves 42 outside
the historical linkage universe. Seven join parents anchored before 2019;
36 remain nonadditive source filings. The remaining 2,347 additive filings,
plus the 2023 Livingston companion inside its parent's observation window,
give the 2,348 canonical constituents. The coverage audit identifies every
excluded source job.

Reweighting uses 595 historical and 310 post-period parents, balancing the
existing log-land-area, residential-FAR, built-FAR and borough moments.
The maximum absolute balance error is $1.42\times10^{-9}$ and the historical
effective sample size is 534.35. The reweighted historical share of parents
totaling exactly 99 units is 0.5261\%, compared with 1.3537\% before the source
correction and 11.9355\% in the post period. This parent-total measure includes
one historical parent whose smaller constituents sum to 99; it differs from
the count of two exact-99 individual filings. Of 499 bootstrap draws, 495
calibrate successfully and four failures are recorded. These are descriptive
comparisons pending the remaining parcel review.

\paragraph{Source corroboration and reviewed boundaries.}
Bedford Square illustrates why an identifier collision needs interpretation.
The archive assigns the same BIN to two of its four documented buildings.
The saved official DOB initial records give all four distinct valid BINs.
The complete additive filing set also matches the reviewed production land
allocation. A common rule permits this combination of evidence while
retaining the archived collision and flagging the corroboration. Bedford is
the only historical parent that qualifies. Its four buildings, 877 units and
177,520.79 square feet remain in the weighting sample.

The historical ownership and description screen also uses 23Q4. Offering-plan
searches cover every historical parent in the canonical period; only plans
with dated pre-snapshot submissions inform that classification. These are
2026 searches of dated records, not an archived 2023 search database. The
supplement adds 652 exact job/address queries and one cross-year Livingston
companion query to the saved August searches. Missing searches remain
different from confirmed zero results. Retrospective July 2026 DOB unit
counts remain labeled comparison fields and do not determine historical units.

\paragraph{Remaining boundary questions.}
The refreshed screen flags 152 of the 905 weighting parents: 120 historical
and 32 post-period. The other 753 pass the current scope checks. These flags
identify questions about coverage and allocation; they are not 152 established
measurement errors. Among the 753 passing parents, 227 candidate land areas
differ from production by more than one square foot. Earlier floor and zoning
allocations still require review before those candidate covariates are adopted.

Of the former 111 queued cases, 103 remain flagged in the weighting sample.
Two of those old cases now belong to one parent, giving 102 distinct current
parents. The other eight have left the weighting sample and have not been
classified as resolved. None of the old queued cases was cleared by this
source correction. The additional 50 flagged parents comprise 29 newly
entering parents and 21 previously supported parents that the revised filing
information sends back for review. Both directions of this reconciliation
are saved as audit tables.

Documentary land reviews also retain their original filing scope. Of 45
reviews in the scope audit, 42 apply to the current weighting parents.
West 48th Street is outside that sample. The two Jamaica/Sutphin reviews
refer to separate filing sets that now form one mechanically linked parent;
their earlier separate allocations do not establish the merged parent's
complete boundary. All 40 production land decisions satisfy their explicit
filing-set checks.

\paragraph{Reproduction.}
The earlier pilot's numerical exhibits are preserved in a dated logbook archive
with input fingerprints. They use the earlier historical source rule. Compiling
the research record does not initiate a new structural fit while boundary
adjudication remains open. A new fit requires the explicit pilot build.

Root \texttt{make} rebuilds the main panels, plots and maps.
\texttt{make dof-site-review} refreshes the parcel checks and reconciles the
old boundary queue with the new sample. The source-rule audit compares
against dated pre-change filing, parent and boundary tables saved in its
\path{code/} directory. The 23Q4 archive is a checksum-pinned production
download, and the new AG query results are fixed received sources. Reviewed
filing roles and land allocations belong to the production manual task.
